%% kapa: exam=1 round=29 year=2018 period=1 form=A subject=real_estate subject_ko=부동산학원론 seq=28 number=108
%% source: Q-Net 감정평가사 기출문제 (cst00309) artlSeq=5202808 file=extracted/1008/5202808/2018년 제29회 감정평가사 1차시험 기출문제/2018년 제29회 감정평가사 1차시험 기출문제/1교시 A형.pdf
%% conversion: kordoc 4.24.51 -> tools/gen1.py (auto); review=proofread-sonnet; flags=none
\begin{question}{108}
부동산개발의 사업타당성분석에 관한 설명으로 옳지 않은 것은?
\begin{choices}
  \item 물리적 타당성분석은 대상 부지의 지형, 지세, 토질과 같은 물리적 요인들이 개발대상 부동산의 건설 및 운영에 적합한지 여부를 분석하는 과정이다.
  \item 법률적 타당성분석은 대상 부지와 관련된 법적 제약조건을 분석해서 대상 부지 내에서 개발 가능한 용도와 개발규모를 판단하는 과정이다.
  \item 경제적 타당성분석은 개발사업에 소요되는 비용, 수익, 시장수요와 공급 등을 분석하는 과정이다.
  \item 민감도분석은 사업타당성분석의 주요 변수들의 초기투입 값을 변화시켰을 때 수익성의 변화를 예측하는 과정이다.
  \item 투자결정분석은 부동산개발에 영향을 미치는 인근 환경요소의 현황과 전망을 분석하는 과정이다.
\end{choices}
\end{question}
\answer{5}
